% TOP_PROBLEM: {"id": 2680, "title": "Kirby Problem 1.21", "category_id": 7, "category": {"id": 7, "name": "topology", "display_name": "Topology", "description": "Properties preserved under continuous deformations.", "slug": "topology", "order_index": 7, "created_at": "2026-07-31T15:26:25.671Z"}, "status": "open"}
% FIRST_POSTED: null
% ENTRY_KIND: statement_only
% Imported from: batch10.tex; UnsolvedMath ID 2680
\documentclass[11pt]{article}
\usepackage[margin=25mm]{geometry}
\usepackage{fontspec}\setmainfont{Latin Modern Roman}
\usepackage{amsmath,amssymb,mathtools}
\usepackage{unicode-math}\setmathfont{Latin Modern Math}
\usepackage{xurl,hyperref}
\hypersetup{colorlinks=true,urlcolor=blue}
\setcounter{secnumdepth}{0}
\setlength{\parindent}{0pt}\setlength{\parskip}{5pt}
\setlength{\emergencystretch}{3em}
\newcommand{\sref}[2]{\hyperlink{source:#1:#2}{[S#2]}}
\newcommand{\cataloguescope}[1]{\paragraph{Recorded target.}#1}
\begin{document}
% UnsolvedMath ID 2680; source catalogue code KP-1.21
\setcounter{section}{2679}
\section{Kirby Problem 1.21}\label{2680}
{\small\sffamily UnsolvedMath ID 2680 · Topology}
\subsection{Definitions and mathematical statement}

\paragraph{Shared notation.}$\mathbb N=\{1,2,\ldots\}$, and $\mathbb Z,\mathbb Q,\mathbb R,\mathbb C$ denote the integers, rationals, reals and complex numbers. Any other conventions are specified locally.


Use the integral hat Heegaard Floer group $\widehat{\mathrm{HF}}(Y)$ and integral hat knot Floer group $\widehat{\mathrm{HFK}}(Y,K)$, with rank meaning the rank of the free abelian part. Knot Floer isomorphisms preserve Maslov and Alexander gradings. A rational homology sphere $Y$ is an L-space when $\operatorname{rank}\widehat{\mathrm{HF}}(Y)=|H_1(Y;\mathbb Z)|$. An L-space knot is a knot in $S^3$ having a positive integral surgery that is an L-space. Let $g(K)$ be the least genus of a Seifert surface. For coprime nonzero integers $p,q$ with $|p|,|q|>1$, $T_{p,q}$ is the torus knot of those parameters; changing one sign gives its mirror. No positivity restriction is imposed on the torus knot in part(c)(ii). An iterated torus knot is obtained from a torus knot by successive cablings. Normalize $\Delta_K(t)$ up to $\pm t^k$.
\paragraph{Statement recorded in the supplied source.}
(a) For every fixed $g\geq1$, are there finitely many L-space knots of genus $g$? (b) Characterize all knots $L$ for which only finitely many isotopy classes $K$ have $\widehat{\mathrm{HFK}}(S^3,K)\cong\widehat{\mathrm{HFK}}(S^3,L)$. (c) Are all of the following true: (i) the knot Floer homology of $T_{2g+1,2}$ detects that knot; (ii) a knot with the knot Floer homology of any $T_{p,q}$ is an iterated torus knot; (iii) an L-space knot all of whose Alexander roots have modulus one is an iterated torus knot? \sref{2680}{1}\sref{2680}{2}
\subsection{Short English statement}
Ask for finiteness at fixed genus or fixed knot Floer homology, and for the three specified detection statements involving torus knots and Alexander roots.
\subsection{Sources}
\begin{description}
\item[\hypertarget{source:2680:1}{S1}] ULAM.AI and the UnsolvedMath Contributors, \emph{UnsolvedMath: A Curated Collection of Open Mathematics Problems}, record 2680 (KP-1.21). CC BY 4.0. \url{https://huggingface.co/datasets/ulamai/UnsolvedMath}.
\item[\hypertarget{source:2680:2}{S2}] R. \.{I}nan\c{c} Baykur, Robion C. Kirby and Daniel Ruberman, K3 -- A New Problem List in Low-Dimensional Topology, author preliminary edition (2026), Problem 1.21 and its remarks. \url{https://math.berkeley.edu/sites/default/files/surv-295-ruberman-watermarked-author-pdf.pdf}.
\end{description}
\label{end:2680}
\end{document}
